% !TeX root = ../papers/nc-p3.tex
\section{可変の跡が3で消える円板における第六経路と四次中心}
\label{sec:rank-fifteen-variable-zero-disk}
\begin{theorem}\label{thm:rank-fifteen-variable-zero-disk}
標数零、剰余標数3の付値を $\nu$ とする。通常な階数15窓の座標が
$0<\nu(y)<3\nu(x)$ を満たすとき、各始点の長さ6の経路の文字別次元は
$(20,20,22,22)$、総次元は84である（文字は以下の規格化で並べる）。
巡回四次中心の文字別次元は $(3,2,0,0)$ であり、各始点で独立な四本を共通に選べる。
\end{theorem}
\begin{proof}
$(U+3)^3=(x^3/y)(U-1)^2$ は整係数のモニック方程式で、還元は $U^3=0$ である。
$U-1$ は単元であり、$K^2=U$, $H^6=y(U-1)^4$ とすると $K,H$ とも正の付値を持つ（$K=0$ も許す）。
$x\ne0$ では $H^2=x(U-1)^2/(U+3)$ を満たす符号・根を選ぶ。
$x=0$ では $U=-3$ として任意の六乗根を取れる。基礎体を拡大し、$i^2=-1$ を剰余体に取る。
$z^4-2Kz^2+1=0$ の単根 $z_*=1+i$ を持ち上げると、
\[
 K=(z^2+z^{-2})/2,\qquad P_0=z+z^{-1},\quad Q_0=z-z^{-1},\quad n_0=(z^2-z^{-2})/2 .
\]
$z,P_0,Q_0,n_0$ はいずれも単元である。そこで計算中は $z$ を不定元とし、
$R=\mathbb Z[z]_{(3,z^2+z+2)}$ とする。$R[[H]]$ において
$P_0^2\pm2H$, $Q_0^2\pm2H$ の平方根を定数項 $P_0,Q_0$ で、
\[
 n=\sqrt{\bigl(n_0^2+\sqrt{n_0^4-4H^2}\bigr)/2}
\]
を内側の定数項 $n_0^2$、外側の定数項 $n_0$ で一意に選べる。
$n_0^2=K^2-1$ なので、$n,H/n$ は積 $H$、平方和 $K^2-1$ を持つ。
これらが通常表示の平方根を与え、経路と中心は正則性座標で比較できる。

\path{computations/verify_rank_fifteen_variable_zero_disk.py} は
$\mathbb Q(z)[H]/(H^7)$ で平方根と関係を構成する。各主元については分母と分子を
$z_*$ で還元し、ともに非零であることを検査する。$H$ による除算では、
捨てる定数項の分子が $\mathbb Z[z]$ の零多項式であることを確かめる。
従ってこれらの行操作は $\bar z=z_*$ を満たすどの持上げにも適用できる。
長さ1から5の各段の最大精度損失は $(0,0,0,2,2)$ であり、全始点で経路基底を得た後に
精度3が残る。この基底で四次中心方程式を文字ごとに消去すると、単元主元だけで
残る自由列数が順に $3,2,0,0$ に達する。文字零の後二自由列と文字一の両自由列に
対応する中心式を各始点で評価すると、単元主元だけで商の長さ4の文字別次元は
$(7,8,8,8)$ に達する。四本が作る行はこれ以上の階数を持ち得ない。長さ6では四始点と四文字の各挿入行列の
幅が56で、最大2回の $H$ 除算で主張の余次元に達する。$H\ne0$ なので除算は
特殊化先の分数体で有効である。経路の各文字の次元には一般点からの下界があり、
長さ4と5の経路次元が一定の部分では中心方程式の核にも同様の下界がある。
以上の消去の上界と一致するので結論を得る。
\end{proof}

\begin{corollary}
$D\ne0$ が3の上の正付値を持つとき、$x^3=Dy$ 上の標数零の全通常窓で
各始点の第六経路の次元は84、四次中心の文字別次元は $(3,2,0,0)$ である。
\end{corollary}
\begin{proof}
係数体の付値を点の座標体まで延長する。$\nu(y)>0$ なら上の定理が適用される。
$\nu(y)=0$ なら $\nu(x)>0$ なので定理\ref{thm:rank-fifteen-residue-line}の経路・中心の計算を使う。
$\nu(y)<0$ なら $3\nu(x)=\nu(D)+\nu(y)$ より $\nu(y)<2\nu(x)$ であり、
定理\ref{thm:rank-fifteen-pole-cone}の経路と中心の計算を使う。
\end{proof}
正付値領域全体での全次数の正則性を決めるには、四次中心式を課した商の高次経路が残る。
